% =============================================================================
\clearpage
\subsection{Trajectoires d'innovation et dépendance de sentier}\label{sec:d10-innovation}
% =============================================================================
% DRAFT v0.3 — mai 2026
% Intègre : skill thesis-dimension-writing, AABH 2012, stratification
% compute/transmit/store, development block, reverse salient, protocole §5.7
% pour quality-energy ladder, diversité d'exemples (pas seulement LLM).


\subsubsection{Le champ~: le facteur limitant de la production de changement technique}

\markEVO \markNEO \markSTS
Le \emph{development block} de la numérisation~--- l'ensemble des activités d'innovation complémentaires que \textcite{nuvolari_understanding_2019}, à la suite de \textcite{carlsson_development_1991}, identifie dans la révolution ICT~--- a fonctionné pendant un demi-siècle comme un réseau autocatalytique à quatre composants~: semi-conducteurs, ordinateurs, logiciels, réseaux. Les innovations dans un composant stimulaient les innovations dans les autres~: la miniaturisation des semi-conducteurs rendait les ordinateurs plus puissants, les ordinateurs plus puissants rendaient les logiciels plus ambitieux, les logiciels plus ambitieux exigeaient des réseaux plus rapides, les réseaux plus rapides créaient un marché pour de nouveaux logiciels~--- et le cycle recommençait. L'énergie ne figurait pas dans le bloc~: elle était un input de fonctionnement, pas un composant du système d'innovation.

\textcite{kander_power_2013} est le seul ouvrage cité par Nuvolari qui intègre explicitement l'énergie comme composant d'un \emph{development block}~--- celui de la Révolution industrielle (vapeur-fer-charbon-machines). Nuvolari note que c'est \og \emph{an important point that deserves further reflection}\fg{}. Le champ que cette dimension ouvre est cette réflexion~: que se passe-t-il quand l'énergie entre dans le \emph{development block} de l'ICT, non comme un enabler~--- le charbon qui nourrit la vapeur~--- mais comme une contrainte sur le rythme d'innovation lui-même~?

La question analytique est double. Le facteur limitant de l'innovation se déplace-t-il du talent~--- le nombre et la qualité des chercheurs, qui contraignait le bloc pendant les phases~1 à~3~--- vers l'énergie~--- la puissance électrique disponible, raccordée et contractualisée~? Et l'orientation du changement technique est-elle verrouillée dans une direction dont le profil énergétique s'alourdit~?

Pour répondre, trois traditions sont mobilisées. \EVO{\textcite{dosi_technological_1982} montre que l'innovation suit des trajectoires à l'intérieur de paradigmes qui définissent les problèmes pertinents et les directions du progrès.} \textcite{acemoglu_directed_2002} formalise la direction du changement technique comme choix endogène aux prix relatifs et à la taille des marchés. \textcite{acemoglu_environment_2012} transposent ce cadre au cas environnemental en distinguant innovation \emph{clean} et innovation \emph{dirty} et en montrant que l'élasticité de substitution entre les deux détermine si une intervention temporaire suffit à rediriger la trajectoire. \STS{\textcite{hughes_networks_1983} introduit le concept de \emph{reverse salient}~--- le composant d'un système technique qui est en retard sur les autres et qui bloque l'ensemble.} Le croisement de ces traditions~--- le \emph{reverse salient} comme diagnostic, le \emph{directed technical change} comme explication, le paradigme comme mécanisme de persistance~--- est une construction propre à cette thèse. Aucune des trois traditions ne la fait seule~: Dosi ne formalise pas la direction, Acemoglu ne pense pas en systèmes techniques, Hughes ne fait pas d'économie de la croissance.


\begin{table}[!ht]
\caption{Deux propriétés du processus d'innovation et leur variation entre régimes.}\label{tab:d10-proprietes}
\centering\footnotesize
\renewcommand{\arraystretch}{1.25}
\begin{tabularx}{\textwidth}{@{}p{2.2cm}p{4.0cm}X@{}}
\toprule
\textbf{Propriété} & \textbf{Question} & \textbf{Ce qui change} \\
\midrule
Direction & Vers quel facteur l'effort d'innovation s'oriente-t-il~? & Du \emph{skill-biased} au \emph{compute-intensive} dans la couche AI~; pas de changement dans les couches réseau et stockage, qui restent en régime d'efficience \\
Persistance & La direction peut-elle changer sous l'effet d'un signal de prix ou de politique~? & Dépend de l'élasticité de substitution entre innovation compute-intensive et innovation energy-saving, et des rendements d'adoption accumulés dans la trajectoire en place \\
\bottomrule
\end{tabularx}
\renewcommand{\arraystretch}{1.0}
\end{table}


\subsubsection{Ce que le matériau révèle}

L'analyse historique fait apparaître trois régularités sur le processus d'innovation dans les techniques de l'information.

La première est que la direction de l'innovation est déterminée par le système métrique du paradigme en place, pas par une évaluation objective de l'espace des possibles. Gray ne voit pas la valeur commerciale du téléphone vocal parce que ni la métrique du paradigme télégraphique~--- messages par fil, coût par mot~---, ni le modèle d'affaires de Western Union ne contiennent la case (section~\ref{sec:telephone}). Gray ne manque pas d'intelligence~; il manque d'unités. Le \emph{focusing device} de \textcite{rosenberg_factors_1969} opère ici en négatif~: il détourne l'attention d'une valeur que les métriques en vigueur ne mesurent pas. Le multiplexage télégraphique obéit au même mécanisme~: la course au débit par fil, orientée par la métrique du paradigme, engendre le téléphone comme sous-produit accidentel~--- la direction de la recherche est fixée par la métrique, pas par la valeur du résultat. Edison, formé dans le paradigme télégraphique, transpose ses compétences vers l'électrification (section~\ref{sec:electrification})~: la direction du changement technique suit la mobilité des personnes entre paradigmes, pas un calcul entre facteurs.

La deuxième régularité est la persistance de trajectoires techniques bien au-delà du point où des alternatives existent. Le système sexagésimal survit cinq millénaires. L'architecture de von Neumann structure le calcul depuis 1945. La carte à quatre-vingts colonnes d'IBM contraint les formats de données pendant un demi-siècle. La loi de 1837 qui réglemente le télégraphe optique s'applique au télégraphe électrique (section~\ref{sec:telegraph})~: la persistance réglementaire est un canal de verrouillage distinct du canal technique, qui traverse les paradigmes.

La troisième régularité, la plus directement pertinente pour cette dimension, est que le coût de l'innovation dans les techniques de l'information n'a jamais été mesuré en énergie. \textcite{nordhaus_two_2007} mesure le coût du calcul sur deux siècles sans inclure l'électricité. \textcite{koomey_implications_2011} mesure l'efficacité énergétique par calcul sans la croiser avec le coût par calcul. L'analyse historique a montré (section~\ref{subsec:occultation}) que cette absence n'est pas un oubli~: c'est le produit des catégories d'analyse en vigueur, qui ne contiennent pas la case énergie-matière. Dans la littérature économique sur l'innovation, le diagnostic est le même~: \textcite{romer_endogenous_1990} traite les idées comme un input non-rival produit par du travail de recherche~; \textcite{aghion_model_1992} mesurent le coût de la destruction créatrice en effort de R\&D~; \textcite{bloom_ideas_2020} mesurent les rendements décroissants de la recherche en nombre de chercheurs. L'énergie n'est le facteur limitant dans aucun de ces cadres. C'est le constat que le matériau de la phase contemporaine force à réviser.


\subsubsection{Deux mécanismes causaux}


\paragraph{Premier mécanisme --- La sélection de la direction par l'effet de taille de marché}

\markNEO

\emph{Observation.} L'analyse historique montre que la direction de l'innovation est orientée par l'environnement institutionnel de l'innovateur. Gray est orienté par le monopole de Western Union et la métrique du débit télégraphique. La course au multiplexage crée le téléphone comme sous-produit~: le marché pour le débit est vaste, le marché pour la voix n'existe pas encore. Edison est orienté par les compétences accumulées dans le paradigme télégraphique~: il transpose, pas il optimise. Dans chaque cas, la direction est sélectionnée par la structure de l'industrie~--- son monopole, ses métriques, la taille de ses marchés~--- pas par un calcul rationnel entre facteurs.

\emph{Cadre théorique.} \textcite{acemoglu_directed_2002} formalise cette sélection. Deux forces concurrentes déterminent la direction~: l'effet de prix, qui oriente l'innovation vers le facteur rare parce que l'économiser est rentable, et l'effet de taille de marché, qui oriente l'innovation vers le facteur abondant parce que le marché pour les innovations complémentaires est plus vaste. L'élasticité de substitution~$\varepsilon$ entre les deux directions tranche~: si $\varepsilon > 1$, la taille du marché domine.

\textcite{acemoglu_environment_2012} transposent ce cadre au cas environnemental. Dans un modèle à deux secteurs~--- \emph{clean} et \emph{dirty}~--- avec innovation dirigée, ils montrent que lorsque les inputs sont suffisamment substituables ($\varepsilon > 1$), une intervention temporaire~--- taxe carbone et subvention à la R\&D propre~--- suffit à rediriger la trajectoire. Le mécanisme est que les rendements d'adoption~--- innover dans un secteur rend l'innovation future dans ce secteur plus rentable, ce qu'ils appellent le \og \emph{building on the shoulders of giants}\fg{}~--- travaillent en faveur du secteur qui a pris l'avance une fois le basculement amorcé. Quand les deux directions sont complémentaires ($\varepsilon < 1$), une intervention permanente est nécessaire. \textcite{popp_induced_2002} fournit l'évidence empirique de l'effet de prix~: sur brevets américains de 1970 à 1994, l'élasticité prix-brevets dans les technologies \emph{energy-saving} est de 0,35 à long terme. \textcite{aghion_carbon_2016} confirment sur données firme dans l'automobile (80~pays, plusieurs décennies)~: les firmes innovent davantage en technologies propres quand elles font face à des prix de carburant plus élevés, et la \emph{path dependence} opère au niveau de la firme~--- l'histoire propre d'innovation prédit la direction future.

\emph{Ce qui est nouveau.} Le cadre AABH traite l'énergie comme un \emph{output} du secteur de production~: le secteur \emph{dirty} produit des biens et émet des polluants. L'innovation y est produite par des chercheurs~; l'énergie n'entre pas dans la fonction de production d'innovation. Ce qui change avec les modèles de fondation à grande échelle, à partir de 2020, est que l'énergie devient un \emph{input} du processus d'innovation lui-même. \textcite{strubell_energy_2019} quantifient le phénomène les premiers~: l'entraînement d'un modèle NLP avec recherche d'architecture neuronale émet environ 284~tonnes de CO$_2$. \textcite{patterson_carbon_2021} raffinent les estimations pour les grands modèles et montrent que le choix du datacenter et du processeur peut réduire l'empreinte d'un facteur~100 à~1\,000, mais que le coût de base croît avec la taille du modèle. Le phénomène ne se limite pas aux modèles de langage~: le calcul haute performance pour la simulation scientifique, le rendu 3D, l'entraînement de modèles de vision ou de modèles de diffusion mobilisent le même type d'infrastructure et la même couche compute AI-intensive.

La transposition du cadre AABH au numérique produit un résultat asymétrique. L'effet de taille de marché pousse vers le \emph{compute-intensive}~: le marché pour les capacités d'IA~--- y compris l'IA prédictive qui migre vers le \emph{fine-tuning} de modèles de fondation, le cloud B2B qui se transforme en plateforme d'inférence, les applications sectorielles qui adoptent les architectures dominantes~--- est global et multi-sectoriel. Le marché pour les innovations d'efficacité~--- modèles frugaux, architectures neuromorphiques, calcul réversible~--- est fragmenté, spécifique, plus petit par segment. L'effet de prix devrait pousser vers l'\emph{energy-saving}~: l'énergie est un coût croissant. Mais cet effet est atténué par les contrats d'achat d'énergie à long terme que la dimension~6 a documentés~: les PPAs stabilisent le coût et affaiblissent le signal de rareté. La résultante est une direction \emph{compute-intensive}, non par irrationalité mais par logique économique endogène.

L'hypothèse propre à cette thèse est que le cadre AABH doit être étendu au cas où l'énergie est un input de l'innovation, pas seulement de la production. Cette extension modifie la structure du problème~: la production d'idées n'est plus un processus à input unique (le travail de recherche de Romer) mais un processus à deux inputs (travail et compute), dont le second est rival et énergivore. La non-rivalité du modèle entraîné~--- résultat de l'innovation~--- masque la rivalité de l'infrastructure qui l'a produit~--- un résultat qui croise la dimension~1. Formaliser cette extension dépasse le périmètre du présent chapitre.

\paragraph{Hypothèse H-D10.1}
\emph{Lorsque le marché pour les capacités de calcul est suffisamment vaste relativement au marché pour les innovations d'efficacité énergétique du calcul, et lorsque le signal de prix de l'énergie est atténué par des contrats de long terme, l'innovation s'oriente vers le compute-intensive par effet de taille de marché au sens d'\textcite{acemoglu_directed_2002}. Ce mécanisme opère dans la couche compute AI-intensive~--- pas dans les couches réseau et stockage, qui restent en régime d'efficience croissante. La condition formelle est $\varepsilon > 1$ entre innovation compute-intensive et innovation energy-saving, condition non vérifiée empiriquement pour cette paire. Si $\varepsilon < 1$~--- par exemple si la distillation de modèles frugaux requiert l'entraînement préalable de gros modèles, rendant les deux directions complémentaires~--- le résultat d'\textcite{acemoglu_environment_2012} est que seule une intervention permanente peut rediriger la trajectoire.}


\paragraph{Deuxième mécanisme --- La persistance par co-évolution techno-institutionnelle et reverse salient}

\markEVO \markSTS

\emph{Observation.} La persistance des trajectoires, telle que l'analyse historique la documente, ne se réduit pas au lock-in technique. La carte à quatre-vingts colonnes d'IBM persiste parce que les données accumulées sont dans un format non-portable~--- c'est du lock-in par le stock, pas par le flux. La loi de 1837 persiste parce que le régulateur applique à un paradigme nouveau les catégories de l'ancien~--- c'est du lock-in institutionnel. \textcite{roussilhe_phase_2025} montre que les \og lois\fg{} contemporaines~--- la loi de Moore, les \emph{scaling laws}~--- sont des dispositifs de coordination performatifs qui alignent les anticipations des acteurs sur une feuille de route commune. Tant que l'industrie y croit, elle investit les sommes nécessaires pour les réaliser.

\emph{Cadre théorique.} \EVO{\textcite{dosi_technological_1982} formalise le paradigme technologique comme un modèle de résolution de problèmes qui concentre l'attention des innovateurs sur certains goulots d'étranglement et détourne l'attention des alternatives hors cadre.} \textcite{arthur_competing_1989} formalise la conséquence dynamique~: les rendements croissants d'adoption verrouillent la trajectoire par accumulation d'investissements, de compétences et d'institutions. \STS{\textcite{unruh_understanding_2000} étend ce mécanisme au domaine de l'énergie sous le nom de \emph{carbon lock-in}~: la co-évolution des systèmes techniques et institutionnels verrouille les économies dans les systèmes fossiles.} \textcite{hughes_networks_1983} apporte un outil diagnostique complémentaire~: le \emph{reverse salient} est le composant d'un système technique qui est en retard sur les autres et qui devient le goulot d'étranglement de l'ensemble.

\textcite{aghion_carbon_2016} fournissent l'évidence empirique de la persistance au niveau firme~: sur données de brevets automobiles dans 80~pays, l'histoire propre d'innovation (ratio brevets \emph{clean}/\emph{dirty}) prédit la direction future de la firme au-delà des effets de prix. C'est le \og \emph{building on the shoulders of giants}\fg{} d'AABH mesuré~: chaque innovation dans une direction rend l'innovation dans cette direction plus rentable, ce qui renforce le verrouillage.

\emph{Ce qui est nouveau.} Le paradigme des \emph{scaling laws}~--- où la performance croît comme une loi de puissance du compute d'entraînement \parencite{kaplan_scaling_2020,hoffmann_training_2022}~--- instancie ce verrouillage dans la couche compute AI. Les \emph{frameworks} dominants sont optimisés pour les GPU, les GPU pour les architectures \emph{transformer}, les architectures \emph{transformer} pour le \emph{scaling}. Les compétences sont formées sur cette pile. Les articles mesurent le progrès par la performance sur des \emph{benchmarks} qui ne contiennent pas de métrique énergétique\footnote{\textcite{schwartz_green_2020} proposent la distinction entre \emph{Red~AI} (performance à tout prix~; le compute a été multiplié par environ 300\,000 entre 2012 et 2018) et \emph{Green~AI} (efficience comme critère d'évaluation). L'initiative MLPerf intègre depuis 2022 un volet efficacité. Ces métriques restent marginales dans les classements qui orientent les investissements.}. Le \emph{focusing device} contemporain est le \emph{benchmark}~: il mesure ce que le paradigme définit comme progrès, pas ce que la contrainte physique impose comme coût.

Le verrouillage ne concerne pas l'ICT dans son ensemble. La couche réseau ne suit pas les \emph{scaling laws}~: ses gains d'efficience compensent la croissance du trafic, et son régime d'innovation reste orienté vers l'efficacité. Le stockage classique reste en régime d'efficience croissante. Le cloud conventionnel~--- machines virtuelles, conteneurs, orchestration~--- poursuit sa trajectoire d'amélioration du PUE. Le minage de cryptomonnaies à preuve de travail est énergivore mais n'obéit pas à une \emph{quality ladder}~: c'est une course à l'efficience du hash dans un protocole fixe, pas un processus de destruction créatrice entre échelons de qualité. L'IA prédictive classique~--- crédit scoring, logistique, maintenance~--- tourne sur du compute conventionnel. Le verrouillage est localisé dans la couche qui concentre le capital et la dynamique d'innovation~: le compute AI-intensive.

L'analogie avec le \emph{carbon lock-in} d'Unruh a une structure précise mais une direction inversée. Chez Unruh, c'est le système énergétique fossile qui verrouille les choix technologiques par inertie~: le gaz verrouille la turbine, la turbine verrouille la centrale, la centrale verrouille le réseau. Ici, c'est la trajectoire d'innovation numérique qui verrouille les besoins énergétiques par escalade~: les \emph{scaling laws} exigent plus de compute, le compute exige plus d'énergie, l'énergie exige de l'infrastructure lourde et irréversible (dimension~3), et cette infrastructure une fois construite stabilise la trajectoire qui l'a appelée. L'énergie est le \emph{reverse salient} du \emph{development block} compute AI~--- le composant dont le rythme de développement est en retard sur le reste du bloc et qui contraint l'ensemble.

\paragraph{Hypothèse H-D10.2}
\emph{Le paradigme des \emph{scaling laws} fonctionne comme un \emph{focusing device}~: il définit le progrès en unités de compute et de performance sur \emph{benchmarks}, sans dimension énergétique. Ce cadrage, combiné aux rendements croissants d'adoption (compétences, \emph{frameworks}, écosystème logiciel, investissements complémentaires irréversibles), verrouille la trajectoire d'innovation de la couche compute AI vers le compute-intensive. L'énergie fonctionne comme \emph{reverse salient} au sens de \textcite{hughes_networks_1983}~: le composant du \emph{development block} dont le rythme contraint les autres. Le verrouillage se relâche si l'une de trois conditions tombe~: (i)~les métriques de progrès intègrent la dimension énergétique (par exemple performance par watt comme critère de classement), (ii)~une architecture s'affranchit de la relation loi de puissance entre compute et performance (neuromorphique, calcul réversible), (iii)~les investissements complémentaires deviennent réversibles. Le relâchement n'est pas observable à ce jour dans la couche compute AI~; il l'est dans les couches réseau et stockage, qui ne sont pas verrouillées.}


\subsubsection{Confrontation à l'évidence}


\paragraph{Évidence historique}

Pour H-D10.1, le cas du multiplexage télégraphique instancie l'effet de taille de marché. Le marché du débit est celui des opérateurs télégraphiques~--- vaste, bien structuré. Le marché de la voix n'existe pas. Gray innove dans le débit, pas dans la voix, bien que la voix soit techniquement à sa portée. Le téléphone naît comme sous-produit. Le cas Edison instancie la \emph{path dependence} mesurée par \textcite{aghion_carbon_2016}~: l'histoire propre de l'innovateur (compétences télégraphiques) prédit sa direction future (électrification), indépendamment des prix relatifs.

Pour H-D10.2, trois cas documentent la persistance avec des durées et des canaux différents. Le verrouillage de la carte à quatre-vingts colonnes (un demi-siècle) opère par le stock de données accumulées dans un format non-portable~: le coût de migration est dans les données, pas dans la machine. La persistance réglementaire (loi de 1837 appliquée au télégraphe électrique) opère par un canal institutionnel distinct du canal technique~: le régulateur projette les catégories de l'ancien paradigme sur le nouveau. Le \emph{focusing device} inversé de Gray opère par un canal cognitif~: le paradigme rend invisible la valeur d'innovations d'un autre registre. Les trois canaux~--- stock, institution, cognition~--- se retrouvent dans le paradigme contemporain des \emph{scaling laws}~: le stock de compétences et de \emph{frameworks}, la structure institutionnelle des classements et du financement, le cadrage cognitif par les \emph{benchmarks}.


\paragraph{Évidence contemporaine}

Pour H-D10.1, la direction du changement technique se lit dans l'allocation des ressources. Le compute d'entraînement des modèles \emph{frontier} a crû d'un facteur supérieur à dix mille entre 2012 et 2024, selon les données d'Epoch~AI\footnote{Epoch~AI reconstruit le compute d'entraînement à partir de publications et de spécifications techniques. La méthodologie est publique~; les estimations pour les modèles propriétaires post-2022 portent une incertitude documentée par les auteurs. C'est la source la plus systématique~; aucune alternative \emph{peer-reviewed} ne couvre le même périmètre.}. Le coût d'un modèle \emph{frontier} dépasse plusieurs centaines de millions de dollars en 2024--2025 \parencite{cottier_rising_2025}. Parallèlement, les investissements en efficacité~--- distillation, quantification, \emph{mixture-of-experts}, refroidissement liquide, processeurs spécialisés~--- existent mais représentent un flux marginal par rapport au \emph{scaling}. L'asymétrie est cohérente avec H-D10.1~: l'effet de taille de marché domine un effet de prix émoussé.

L'objection de l'inférence mérite une réponse. \textcite{patterson_carbon_2021} rapportent que l'inférence représente environ trois cinquièmes de la consommation ML totale chez Google, et que le coût unitaire par inférence diminue rapidement. Le matériau de la dimension est le processus d'innovation~--- le passage d'un échelon de qualité au suivant, c'est-à-dire l'entraînement du prochain modèle ou du prochain système. L'inférence est le déploiement du résultat, pas sa production. Que le déploiement devienne plus efficient est un résultat de la dimension~2 (rendements) et de la dimension~5 (effet rebond). Que la production de l'échelon suivant coûte de plus en plus d'énergie est le fait propre de cette dimension. Les deux sont vrais simultanément.

Pour H-D10.2, le verrouillage est observable dans la structure de l'écosystème. \textcite{schwartz_green_2020} documentent que le compute requis pour la recherche en \emph{deep learning} a été multiplié par environ 300\,000 entre 2012 et 2018, et que les progrès sur les \emph{benchmarks} sont systématiquement corrélés à l'augmentation du compute, pas à l'amélioration de l'efficience algorithmique. Six ans après leur proposition de faire de l'efficience un critère d'évaluation, les classements qui orientent les investissements n'ont pas changé de métrique.

Le contradicteur contemporain le plus direct est DeepSeek, qui a montré en 2024--2025 qu'un modèle \emph{frontier} peut être entraîné avec un budget de compute substantiellement inférieur au standard américain. Ce cas ne réfute pas H-D10.2~--- il en précise le domaine. Le verrouillage décrit par H-D10.2 est un verrouillage de la trajectoire d'allocation du capital et de la R\&D dominante~; il n'exclut pas que des trajectoires alternatives émergent sous d'autres contraintes institutionnelles (dans ce cas, la restriction d'accès aux puces avancées). La question que DeepSeek pose à H-D10.1 est plus profonde~: si $\varepsilon > 1$ (les trajectoires frugales substituent les trajectoires intensives), le verrouillage est réversible et le paradigme peut basculer sans intervention permanente. Si $\varepsilon < 1$ (les modèles frugaux sont obtenus par distillation depuis des gros modèles, donc complémentaires), les deux coexistent sans que l'une déplace l'autre. C'est une question empirique que la modélisation devra paramétrer.


\subsubsection{Fait stylisé}

\textcite{bloom_ideas_2020} établissent que la productivité de la recherche décline dans chaque domaine examiné~: pour les semi-conducteurs, le nombre effectif de chercheurs a été multiplié par dix-huit entre 1971 et 2017 pour maintenir le doublement de la densité tous les deux ans. Les \emph{scaling laws} \parencite{kaplan_scaling_2020,hoffmann_training_2022} montrent que dans le paradigme IA, les rendements décroissants s'expriment en compute~: réduire la \emph{loss} de moitié requiert un ordre de grandeur de compute supplémentaire, qui se traduit directement en consommation électrique. Le croisement de ces deux résultats~--- rendements décroissants de la recherche, matérialisation du coût en énergie~--- fonde le fait stylisé.

Avant de le formuler, un concept doit être posé. En combinant la \emph{quality ladder} de \textcite{aghion_model_1992}~--- où chaque échelon rend le précédent obsolète~--- avec les rendements décroissants de \textcite{bloom_ideas_2020} et la matérialisation énergétique du compute, cette thèse propose le terme de \emph{quality-energy ladder} pour désigner le mécanisme par lequel chaque échelon de qualité dans la couche compute AI requiert un saut disproportionné en consommation d'énergie.

Le terme appelle quatre précisions. Premièrement, il est en tension avec son ascendant~: la \emph{quality ladder} d'Aghion et Howitt mesure le coût de l'innovation en effort de recherche~--- un input immatériel. La \emph{quality-energy ladder} le matérialise en énergie. La tension est le signal que la propriété que la \emph{quality ladder} standard traite comme non-rivale~--- l'idée~--- est devenue rivale dans ses conditions de production. Deuxièmement, quatre termes alternatifs ont été considérés et écartés~: \emph{energy-intensive innovation} (ne capture pas la structure en échelons ni la destruction créatrice), \emph{compute ladder} (ne capture pas le lien avec l'énergie), \emph{scaling cost curve} (descriptif, sans dimension de destruction créatrice), \emph{energy return on innovation} (capture le ratio, pas la dynamique séquentielle). Troisièmement, ce que le terme désigne positivement est la conjonction de trois propriétés~: une séquence d'échelons où chaque génération rend la précédente obsolète, un coût croissant en énergie par échelon, et une destruction créatrice entre échelons. Les trois sont établies séparément~; leur conjonction dans un même processus est le phénomène que le terme désigne. Quatrièmement, c'est une proposition de cette thèse, pas un concept établi. Une recherche dans la littérature n'a identifié aucun usage préalable. Les précurseurs les plus proches sont \textcite{schwartz_green_2020} (explosion du compute sans structure en échelons), \textcite{strubell_energy_2019} (coût carbone de l'entraînement sans lien avec les rendements décroissants), \textcite{patterson_carbon_2021} (empreinte de modèles spécifiques sans question de trajectoire).

La \emph{quality-energy ladder} est une propriété de la couche compute AI, pas de l'ICT en général. Les couches réseau et stockage ne montent pas cette échelle~--- elles restent en régime d'efficience croissante. Le cloud conventionnel, l'IA prédictive classique, les infrastructures de communication obéissent à d'autres dynamiques. Le diagnostic de \textcite{gordon_does_2000}~--- les TIC ne sont pas une GPT au même titre que l'électricité~--- se repositionne à la lumière de cette stratification. Gordon avait peut-être raison tant que le paradigme Dennard tenait~: l'ICT améliorait l'efficacité sans imposer de contrainte matérielle. Ce qui change n'est pas que Gordon avait tort, c'est que la condition de validité de la thèse adverse~--- le facteur clé bon marché dont le prix baisse continuellement \parencite{bresnahan_general_1995}, hypothèse de \textcite{perez_technological_2002}~--- cesse de tenir pour la couche qui concentre le capital.

\textcite{brynjolfsson_artificial_2019} proposent l'hypothèse du \emph{J-curve}~: les GPT nécessitent des co-investissements massifs avant de produire des gains de productivité. L'argument est cohérent pour le capital organisationnel, dont \textcite{bresnahan_information_2002} estiment le délai à environ dix~ans. La \emph{quality-energy ladder} en étend la portée~: le capital complémentaire en Phase~4 est infrastructurel-énergétique, et le délai passe à quinze-vingt~ans~--- construire un datacenter à l'échelle du gigawatt et raccorder la capacité électrique prend une décennie de plus que réorganiser un bureau.

Le fait stylisé se formule ainsi. Le mécanisme d'improvement de la GPT-ICT se stratifie en Phase~4. La couche compute AI-intensive bascule dans un régime où le coût énergétique de l'innovation est croissant~--- chaque échelon de la \emph{quality-energy ladder} requiert un saut disproportionné en énergie~--- et où la direction de l'innovation est verrouillée vers le \emph{compute-intensive} par l'effet de taille de marché et la co-évolution techno-institutionnelle. L'énergie fonctionne comme \emph{reverse salient} du \emph{development block} de cette couche~: le composant dont le rythme de développement contraint le rythme de l'innovation. Les couches réseau, stockage et compute conventionnel restent en régime d'efficience croissante et ne sont pas affectées par ce mécanisme. La Phase~4 est définie par le fait que la couche qui a basculé est celle qui concentre le capital et la dynamique d'innovation~--- elle tire l'ensemble du système malgré sa localisation dans un segment spécifique. La persistance de ce régime est conditionnelle au maintien du paradigme des \emph{scaling laws}~: si les gains d'efficience (distillation, \emph{mixture-of-experts}, architectures alternatives) compensent la croissance de la demande, la couche compute AI re-basculerait en régime d'efficience et la Phase~4 serait un épisode transitoire. C'est le paramètre de scénario que la modélisation devra instruire.


\begin{table}[!ht]
\caption{Les deux mécanismes de la dimension et leurs hypothèses.}\label{tab:d10-mecanismes}
\centering\footnotesize
\renewcommand{\arraystretch}{1.25}
\begin{tabularx}{\textwidth}{@{}p{2.2cm}p{2.8cm}p{3.6cm}X@{}}
\toprule
\textbf{Mécanisme} & \textbf{Propriétés} & \textbf{Conditions} & \textbf{Hypothèse} \\
\midrule
M1. Sélection de direction & Direction & Marché compute-intensif $\gg$ marché efficacité \emph{et} signal de prix atténué (PPAs) \emph{et} $\varepsilon > 1$ entre les deux directions & H-D10.1 \\
M2. Persistance et reverse salient & Persistance & Métriques sans dimension énergétique \emph{et} rendements croissants d'adoption \emph{et} investissements irréversibles & H-D10.2 \\
\bottomrule
\end{tabularx}
\renewcommand{\arraystretch}{1.0}
\end{table}


\textcolor{teal}{Question ouverte pour la modélisation~: comment représenter dans IMACLIM un changement de paradigme~--- la rupture des \emph{scaling laws}~--- qui est une discontinuité, pas un paramètre continu~? Le module numérique devra distinguer au moins deux régimes (Koomey et post-Koomey) pour la couche compute AI, avec un seuil de basculement. La forme de ce seuil~--- endogène à l'investissement, ou exogène comme hypothèse de scénario~--- est une décision de design du modèle.}

\textcolor{teal}{Refs non intégrées~: David 1990 (\emph{AER})~--- analogie dynamo/ordinateur, déjà mobilisée en D4 et dans la méthode du Ch.2. Klein \& \c{S}ener 2022 (\emph{RED})~--- innovation, diffusion et croissance endogène. Caselli 1999 (\emph{AER})~--- révolutions technologiques et inégalités, croisement D10/D4. Helpman \& Trajtenberg 1998~--- \emph{sow and reap}. Ces références pourraient enrichir le champ ou la confrontation mais ne font pas, en l'état, un travail analytique précis dans l'argument.}
